\documentclass[a4paper,12pt]{article}

\usepackage{alltt, fancyvrb, url}
\usepackage{graphicx}
\usepackage[utf8]{inputenc}
\usepackage{float}
\usepackage{hyperref}
\usepackage[italian]{babel}
\usepackage[table]{xcolor}
\usepackage{array}
\usepackage{multirow}
\usepackage{titlesec}
\usepackage{listings}
\usepackage[italian]{cleveref}
\usepackage{acronym}
\usepackage{subcaption}

\renewcommand{\thesection}{\arabic{section}}

\sloppy
\setlength{\parindent}{0pt}

\title{\textbf{Final Project Report}}
\author{}
\date{}

\begin{document}
\maketitle
\tableofcontents

\newpage

\section{Procedura di accettazione}
Il cliente effettua il collaudo della soluzione consegnata, verificando il rispetto dei criteri di accettazione
stabiliti ed approvati in fase di planning, oltre alle CoS stabilite all'interno del POS. In caso positivo, il cliente
certifica l'accettazione della soluzione consegnata.

\section{Modalità di installazione}
L'installazione della soluzione avviene seguendo un \textit{Phased approach}, quindi verrà installata l'ultima parte
della soluzione, in quanto le prime parti sono già state consegnate ed installate in precedenza.

\section{Audit post-implementazione}
Si svolge l'audit post-implementazione, che consiste in un'analisi della soluzione consegnata e del processo di
sviluppo, in modo da effettuare una retrospettiva per migliorare in progetti futuri. \\
Un esempio di audit post-implementazione è il seguente:
\begin{itemize}
    \item gli obiettivi sono stati raggiunti e il deliverable era conforme con le aspettative del cliente;
    \item il progetto è stato completato rispettando i limiti di budget e tempo;
    \item il committente si è mostrato soddisfatto del risultato;
    \item Kanban non era stato utilizzato in progetti passati, dunque è stata un'opportunità per migliorare la
    conoscenza della metodologia. In particolare, sono stati notati i vantaggi di flessibilità legati all'assenza di
    durate fisse dei cicli di sviluppo caratterizzanti Scrum (metodologia più utilizzata in passato). Inoltre, il team
    ha accolto positivamente questo nuovo approccio, in quanto ha permesso di lavorare in maniera più autonoma, e con
    un maggior senso di responsabilità grazie all'orizzontalità del team.
\end{itemize}

\section{Opportunità future}
L'applicazione prodotta offre diverse opportunità di utilizzo futuro, tra cui:
\begin{itemize}
    \item riutilizzo della logica di registrazione e login degli utenti;
    \item riutilizzo dei meccanismi di sicurezza e protezione dei dati utilizzati;
    \item riutilizzo delle configurazioni dei container e deployment dei vari microservizi.
\end{itemize}
Inoltre, l'aver usato per la prima volta l'approccio Kanban ha permesso di produrre della documentazione sulla base
della quale sarà possibile ricavare template per progetti futuri.

\end{document}